% Reemplazo de la tabla y el párrafo de lectura de la subsección "Comparación con la solución propuesta"
% (sections/contextoProyecto.tex). Misma label: tab:comparacion-soluciones.
La Tabla~\ref{tab:comparacion-soluciones} contrasta a las seis plataformas con Tukit en trece criterios,
repartidos en tres grupos: cómo entra la lista, cómo se arma y se recomienda el kit, y qué tan madura es
la operación comercial. Se incluyen criterios donde Tukit pierde, para que la comparación sea equilibrada.

\begin{table}[htbp]
\centering\footnotesize
\setlength{\tabcolsep}{2pt}
\renewcommand{\arraystretch}{1.2}
\newcommand{\nc}{\textsf{\scriptsize n/c}}
\newcommand{\enc}[2]{{\scriptsize\begin{tabular}[b]{@{}c@{}}#1\\#2\end{tabular}}}
\newcommand{\grupo}[1]{\multicolumn{8}{@{}l}{\cellcolor{gray!8}\textbf{#1}} \\}
\caption{Comparación de Tukit con las soluciones existentes en trece criterios.}
\label{tab:comparacion-soluciones}
\begin{tabular}{@{}>{\raggedright\arraybackslash}p{4.9cm}
  *{6}{>{\centering\arraybackslash}p{1.1cm}}
  >{\columncolor{gray!15}\centering\arraybackslash}p{1.1cm}@{}}
\toprule
\enc{}{\textbf{Criterio}} &
\enc{Teacher-}{Lists} &
\enc{1st Day}{Supplies} &
\enc{Pana-}{mericana} &
\enc{Grupo}{Éxito} &
\enc{Papelería}{Primavera} &
\enc{Mercado}{Libre} &
\enc{\textbf{Tukit}}{} \\
\midrule
\grupo{A. Entrada de la lista}
Parte de la lista oficial del colegio
  & \cumple & \cumple & \nocumple & \nocumple & \nocumple & \nocumple & \cumple \\
Lee listas en foto, PDF o texto
  & \nocumple & \nocumple & \nocumple & \nocumple & \nocumple & \nocumple & \cumple \\
Integración con el colegio o el docente
  & \cumple & \cumple & \nocumple & \nocumple & \nocumple & \nocumple & \nocumple \\
\midrule
\grupo{B. Armado y recomendación}
Asocia cada renglón con un producto del catálogo
  & \parcial & \parcial & \nocumple & \nocumple & \nocumple & \nocumple & \cumple \\
Arma el kit completo
  & \parcial & \cumple & \nocumple & \parcial & \parcial & \parcial & \cumple \\
Combina productos de varias tiendas
  & \parcial & \nocumple & \nocumple & \nocumple & \nocumple & \parcial & \cumple \\
Ofrece alternativas según presupuesto y preferencias
  & \nocumple & \nocumple & \nocumple & \parcial & \nocumple & \parcial & \cumple \\
Explica el kit y revisa los casos dudosos
  & \nc & \nc & \nc & \nc & \nc & \nc & \cumple \\
\midrule
\grupo{C. Operación comercial y madurez}
Tamaño y profundidad del catálogo
  & \parcial & \nc & \cumple & \cumple & \parcial & \cumple & \nocumple \\
Inventario propio y logística de entrega
  & \nocumple & \nc & \cumple & \cumple & \nc & \cumple & \nocumple \\
Aplicación móvil nativa
  & \nc & \nc & \nc & \nc & \nc & \nc & \nocumple \\
Pago real en producción
  & \nocumple & \nc & \cumple & \cumple & \nc & \cumple & \nocumple \\
Madurez y base de usuarios
  & \cumple & \nc & \cumple & \cumple & \nc & \cumple & \nocumple \\
\bottomrule
\end{tabular}

\vspace{3pt}
\begin{minipage}{0.97\textwidth}
{\footnotesize \cumple\ Cumple \quad \parcial\ Cumple parcialmente \quad \nocumple\ No cumple \quad
\nc\ No consta en las fuentes públicas revisadas.
Fuentes: sitios públicos de cada solución, consultados en noviembre de 2026
\citep{teacherlists_retail_2026,firstday_faqs_2026,panamericana_cotizador_2025,exito_listas_2025,primavera_armatulista_2026,mercadolibre_kits_2025}.
TeacherLists y 1st Day no venden en Colombia. Una celda \parcial\ indica que la función existe con restricciones
(por ejemplo, 1st Day solo atiende a colegios aliados). Donde la página no permitió confirmar una función,
se anota n/c y no se supone.}
\end{minipage}
\end{table}

Tukit gana en los criterios de entrada y de armado: es la única de las siete que lee la lista en foto, PDF o
texto, la asocia con el catálogo, combina tiendas y ofrece varias alternativas por presupuesto, de modo que el
acudiente no traduce renglón por renglón. En cambio, las plataformas comerciales lo superan en catálogo,
inventario, entrega, pago real y base de usuarios, porque venden y despachan, mientras que Tukit no almacena,
no entrega y, hasta ahora, no cobra. TeacherLists y 1st Day, además, tienen un vínculo con colegios y docentes
que Tukit aún no construye. Esa diferencia es la oportunidad: las soluciones con operación madura no parten de la
lista oficial, y las que sí parten de ella no operan en Colombia ni leen formatos libres. Por eso Tukit puede
apoyarse en la operación de las tiendas existentes, y su aporte queda en convertir cualquier lista escolar en
kits completos y auditables, que es la brecha de la que parte el levantamiento de requerimientos del
Capítulo~IV.\par
